\documentclass[11pt]{article}
\usepackage[final]{acl}
\usepackage{times}
\usepackage{latexsym}
\usepackage[T1]{fontenc}
\usepackage[utf8]{inputenc}
\usepackage{microtype}
\usepackage{amsmath}
\usepackage{booktabs}
\title{Efficient CV--Job Fit Classification through Distillation and Quantization}
\author{Rawan Khojah \quad Hams Aljohani\\Effat University\\CS4085 Deep Learning}
\begin{document}
\maketitle
\section{Motivation and related work}
HireMind is our Senior recruitment platform. We propose investigating a compact CV--job fit classifier that could later contribute evidence to its scoring workflow. The DL contribution concerns model training and efficiency; platform integration remains future work. It does not replace skill verification or recruiter decisions. A compact local classifier could reduce reliance on external API calls for matching and keep this inference step within an organization. We will measure storage and runtime; local execution alone does not guarantee privacy or lower total cost.

\citet{gan} provide recruitment-screening context. BERT supplies the architecture foundation \citep{bert}; \citet{turc} motivate pretrained compact students. \citet{hinton} explain teacher-guided distillation, and Q8BERT \citep{q8} studies quantization in BERT. Our question is whether a small classifier benefits from distillation on limited CV--job pairs and retains that gain after simpler post-training quantization. These studies do not establish whether the distillation gain survives embedding-inclusive post-training quantization on our CV--job benchmark. We test that combination using a matched student architecture and held-out job descriptions, without claiming a new algorithm.

\section{Question and proposed approach}
\textbf{Research question:} Can teacher-guided training improve compact CV--job fit classification, and can quantization reduce storage while preserving the improvement?

Input is a CV and job description; output is \emph{No Fit}, \emph{Potential Fit} or \emph{Good Fit}. The student is Google's pretrained BERT-Tiny checkpoint, \texttt{bert\_uncased\_L-2\_H-128\_A-2}, with a three-class head. A separately fine-tuned \texttt{bert-base-uncased} supplies teacher logits.

\textbf{M0:} Student trained with hard-label cross-entropy. \textbf{M1:} The same student architecture and matched initialization, trained with
\begin{equation}
L=0.5L_{CE}+0.5T^2D_{KL}(p_T\Vert p_S),
\end{equation}
where teacher and student probabilities use temperature $T$. \textbf{M2:} A CPU copy of M1 with dynamic INT8 Linear modules and 8-bit weight-only Embedding modules. Layer normalization and unsupported operations remain FP32. Embeddings use a dedicated weight-only configuration. M2 requires no additional training; unlike Q8BERT, we use post-training quantization.

The training strategy is Part 1; quantizing its resulting model is Part 2. The teacher is additional support, not M0. All models use identical text-pair preprocessing and a maximum of 256 tokens. Truncation frequency will be reported.

\section{Hypotheses}
Let $F_i$ denote mean test macro-F1 on a 0--100 scale and $S_i$ serialized model size.

\textbf{H1:} $F_1-F_0\geq1.0$ percentage point.

\textbf{H2:} $S_2\leq0.40S_1$, $F_1-F_2\leq1.0$ point, and $F_2-F_0\geq0.5(F_1-F_0)$, with $F_1>F_0$.

These targets are exploratory, not pilot-derived. H2 targets at least 60\% less storage: quantizing embeddings covers the dominant weight table, while metadata and remaining FP32 values limit compression. We do not promise fourfold compression or faster inference. If M1 does not improve, the improvement-retention claim fails.

\section{Data and fair experiments}
We use the instructor-approved resume--job-description fit dataset uploaded by cnamuangtoun on Hugging Face.\footnote{\url{https://huggingface.co/datasets/cnamuangtoun/resume-job-description-fit}} Its inspected revision contains 8,000 pairs: 6,241 training and 1,759 test rows, with 643 distinct CV texts and 351 job texts. Labels have counts 4,000/2,000/2,000. Six pairs have conflicting labels; quarantine their twelve rows and remove one remaining duplicate, leaving 7,987 pairs before further near-duplicate screening.

The official split shares 476 CV texts but no exact job texts. Preserve test jobs and reserve about 20\% of training-job groups for validation with seed 42. Exact and whitespace-normalized checks were completed; inspect near duplicates before freezing indices. This evaluates \emph{unseen jobs with potentially familiar CVs}, not wholly unseen candidates. The instructor approved its use for this project. Its original provenance, licence and labeling method remain undocumented. We will audit training labels, avoid redistributing CV text, and limit claims to the supplied benchmark labels.

Use AdamW, effective batch size 16 and at most three epochs. Give M0 and M1 four validation-based hyperparameter trials each, with equal update caps. Tune learning rate for M0; tune learning rate and $T\in\{2,4\}$ for M1. Select checkpoints using validation macro-F1, then repeat selected settings at seeds 42, 43 and 44. Train the teacher on training data only; cache logits for those same examples. Freeze decisions before final test evaluation.

Report macro-F1, accuracy, per-class scores, confusion matrices, logical parameter count, serialized size, MACs/FLOPs under an explicit convention, CPU median/p95 latency and peak memory. Benchmark identical inputs, batch size, thread count and software, with 30 warm-up and 200 timed calls. Separate tokenization time and teacher-training cost. Report paired-seed variability; use job-grouped uncertainty estimates if feasible. Precision reduction need not reduce operation counts.

\section{Feasibility and team plan}
Completed preparation includes full-text paper review and dataset auditing. No model has been trained and no pilot performance is claimed. A Colab pilot will check checkpoint loading, embedding quantization, save/reload, finite predictions, memory, runtime and validation performance on a small subset. Free Colab is the intended starting point, subject to resource availability.

Our two-person team is approved by the instructor. Rawan leads data auditing and measurement; Hams leads training and export; both verify results and write, recording actual contributions. We selected the Technical Report track. GitHub and course-template Overleaf projects have been created and shared with the instructor. Proposed course-week allocation, subject to calendar confirmation: Week 10, labels and grouped split; Week 11, pilot, teacher and M0; Week 12, M1; Week 13, M2 and profiling; Week 14, final evaluation; Week 15, technical report, presentation and poster.

\section{Risks and fallback}
Undocumented data provenance and labels limit external validity; seek clarification and change data only with instructor agreement if necessary. If teacher quality is weak, report negative transfer or revise to class-weighted training before opening the test set. If embedding quantization fails compatibility checks, retain Linear-only quantization as a documented fallback and revise H2 before test evaluation; report any quality loss honestly. Reduce batch size with gradient accumulation for memory limits. Treat truncation and label ambiguity as evaluation limitations. Predictions will not support autonomous hiring or claims of demographic fairness.

\section{AI use}
ChatGPT and Claude assisted literature analysis, dataset research and auditing, experiment planning, workspace setup and proposal drafting. The authors are responsible for checking the content and references and for the final submission. No experimental results were generated or claimed.

\clearpage
\begin{thebibliography}{5}
\bibitem[Gan et al.(2024)]{gan} Chengguang Gan, Qinghao Zhang, and Tatsunori Mori. 2024. Application of LLM Agents in Recruitment: A Novel Framework for Resume Screening. \emph{arXiv:2401.08315v2}. \url{https://arxiv.org/abs/2401.08315v2}.
\bibitem[Devlin et al.(2019)]{bert} Jacob Devlin, Ming-Wei Chang, Kenton Lee, and Kristina Toutanova. 2019. BERT: Pre-training of Deep Bidirectional Transformers for Language Understanding. \emph{NAACL-HLT}, pages 4171--4186. \url{https://aclanthology.org/N19-1423/}.
\bibitem[Turc et al.(2019)]{turc} Iulia Turc, Ming-Wei Chang, Kenton Lee, and Kristina Toutanova. 2019. Well-Read Students Learn Better: On the Importance of Pre-training Compact Models. \emph{arXiv:1908.08962v2}. \url{https://arxiv.org/abs/1908.08962v2}.
\bibitem[Hinton et al.(2015)]{hinton} Geoffrey Hinton, Oriol Vinyals, and Jeff Dean. 2015. Distilling the Knowledge in a Neural Network. \emph{arXiv:1503.02531}. \url{https://arxiv.org/abs/1503.02531}.
\bibitem[Zafrir et al.(2019)]{q8} Ofir Zafrir, Guy Boudoukh, Peter Izsak, and Moshe Wasserblat. 2019. Q8BERT: Quantized 8Bit BERT. \emph{EMC2 workshop}. DOI: 10.1109/EMC2-NIPS53020.2019.00016. \url{https://arxiv.org/abs/1910.06188v2}.
\end{thebibliography}
% After student revision and the required NEW-CHAT review, create ai_review.tex
% containing the final review report (not an invented score).
\IfFileExists{ai_review.tex}{\clearpage\appendix\section{AI review report}\noindent\textit{Review generated by ChatGPT on 18 October 2026 in a new conversation, using the course AI Reviewer prompt and the complete proposal text. Web-citation labels inserted by the chat interface were removed; the review content is otherwise reproduced unchanged.}

\subsection*{CS4085 Deep Learning — Formative Proposal Review}
Project: HireMind: Efficient CV–Job Fit Classification through Distillation and Quantization

Review standard: Strict, rubric-based assessment of the proposal as written, not an official course grade. I am evaluating research design, technical correctness, reproducibility, feasibility, and the strength of the evidence promised by the plan. I am not requiring full model training at the proposal stage.

\subsubsection*{1. Overall assessment}
Provisional formative score

\subsection*{83/100}
READY TO SUBMIT*

*Under the stated numerical threshold, provided the technical and experimental details below are corrected or verified before submission.

The proposal is unusually well specified for a project that has not yet trained a model. It defines all three models, connects the two contributions, sets numerical hypotheses, distinguishes validation from test data, and plans reproducible measurements.

The main weakness is not missing structure. It is the gap between the intended quantization experiment and what the proposed implementation and pilot can reliably establish. The literature gap also needs a more careful qualification, and the experimental design needs stronger safeguards against candidate-level and dataset-level leakage.

My strict-reviewer verdict: a strong proposal with several important technical and methodological details to resolve. It should not be treated as a completed research result. The absence of trained models is acceptable at this stage, and you explicitly disclose that no results exist.

\subsubsection*{2. Category breakdown}
\begin{small}\begin{tabular}{p{0.25\linewidth}p{0.12\linewidth}p{0.58\linewidth}}\toprule
\textbf{Rubric category} & \textbf{Score} & \textbf{Reviewer assessment} \\\midrule
Motivation & 9/10 & Practical motivation, bounded claims, and a clear efficiency rationale. \\
Related work (at least 5 papers) & 8/10 & Five sources are cited, but the synthesis and novelty claim need stronger verification. \\
Research question and gap & 9/10 & Specific, testable question and a clear proposed gap. \\
Deep learning contribution & 13/15 & M0/M1 are well defined; teacher details and the interpretation of the gain need tightening. \\
Efficiency contribution & 11/15 & M2 is specified, but quantization support, implementation details, and the compression target need validation. \\
Two hypotheses & 9/10 & Numerical, connected, and appropriately conditional. \\
Dataset and fair experiments & 7/10 & Strong split discipline, but repeated-CV overlap, label validity, and the tuning protocol deserve more scrutiny. \\
Feasibility and team plan & 8/10 & Sensible pilot and division of work; the proposed quantization path remains a technical risk. \\
Risks and fallback & 4/5 & Honest fallbacks; the Linear-only fallback changes the efficiency method and cannot establish the original H2. \\
Writing, format, AI disclosure & 5/5 & Clear structure, measurable claims, and transparent AI-use disclosure. \\
Total & 83/100 & Strong proposal; revisions recommended. \\
\bottomrule\end{tabular}\end{small}\par\medskip
These are formative estimates, not an official grade. The score reflects the written plan, not experimental performance.

\paragraph*{Mandatory-rule audit}
M0, M1, and M2 defined

M0 is the hard-label student baseline, M1 adds teacher-guided distillation, and M2 quantizes M1. Their relationship is explicit.

H1 compares M1 against M0

$F_1-F_0\geq1.0$ macro-F1 point. This directly tests the deep learning contribution.

H2 compares M2 against M1 and assesses M2 against M0

H2 tests size reduction, the loss relative to M1, and retention of at least half the M1-versus-M0 gain. This is a meaningful connected efficiency contribution.

Numerical targets present

The proposal gives explicit thresholds for both hypotheses and a 95\% validation-prediction agreement criterion for the pilot.

No stated test-set tuning

Hyperparameter selection is assigned to validation, and test access is deferred until choices are frozen. This is the correct intended protocol.

One important qualification: these checks confirm that the proposal states the required design. The implementation, data audit, and eventual experiment logs must demonstrate that the design is followed.

\subsubsection*{3. Three main improvements}
\paragraph*{1. Prove that the planned quantization works on the actual model}
This is the most important technical issue. The proposal specifies dynamic INT8 quantization for Linear layers and weight-only INT8 quantization for Embedding layers. Those are not interchangeable operations, and successful conversion of one does not establish that the full Hugging Face model will convert and run correctly.

PyTorch provides dedicated embedding quantization support and configuration requirements. Verify compatibility with your exact model architecture, installed PyTorch version, and save/reload path.

Your existing pilot is a good start. Strengthen its evidence by confirming:

\begin{itemize}
\item The actual embedding weight dtype and representation after conversion, rather than relying only on a successful conversion call.
\item Successful inference after saving and reloading the entire model.
\item The actual serialized-size ratio $S_2/S_1$.
\item Finite outputs and macro-F1 on a small validation subset, alongside the prediction-agreement check.
\item The exact modules quantized and those left in FP32.
\end{itemize}
Why this matters: If embeddings are not quantized as intended, M2 may miss its compression target even if Linear quantization works. A passing conversion test is not the same as a passing efficiency experiment.

\paragraph*{2. Make the literature gap more defensible}
The proposal cites five relevant sources, and the combination of distillation and quantization is a plausible experimental question. However, the claim that no prior work examines the combination is stronger than the evidence provided.

Q8BERT studies quantization-aware training, whereas your proposed M2 uses post-training quantization. That difference helps motivate your experiment, but it does not by itself establish that the overall combination is novel. The compact-model distillation literature also includes work beyond the one result quoted from Turc et al.

Improve the research argument by verifying whether prior work has studied:

\begin{itemize}
\item Distilled compact BERT models followed by post-training quantization.
\item Quantization of embedding weights in compact Transformer models.
\item The interaction between distillation gains and post-training quantization, even on tasks other than recruitment.
\end{itemize}
Then qualify the gap to match what you can substantiate. You do not need to claim a new algorithm: a well-controlled, task-specific empirical study is a reasonable course contribution.

\paragraph*{3. Strengthen the data and evaluation argument}
The split protocol is more thoughtful than a random row-level split, but the proposal reports 476 shared CVs between the official training and test portions. It correctly acknowledges that the experiment therefore measures performance on unseen jobs with potentially familiar CVs.

That is an important limitation, not a trivial detail. A classifier can potentially exploit resume-specific patterns when a CV appears in both training and test data. This may affect how you interpret the measured fit-classification performance.

Before training, ensure the audit clearly documents:

\begin{itemize}
\item Exact CV and job overlap after deduplication and conflict removal.
\item The final train, validation, and test counts after every filtering operation.
\item Whether your intended claim is generalization to unseen jobs, unseen candidates, or both.
\item How the reported results might change if candidate-level separation were used.
\end{itemize}
The existing protocol can support a narrowly scoped study of unseen jobs. It cannot, on its own, establish generalization to unseen candidates. Keep that distinction explicit in your report.

\subsubsection*{4. Major and minor issues}
\paragraph*{Major issues}
These are the items I would expect a strict instructor to challenge.

Major 1The full-model quantization path is not yet demonstrated.

The proposed API and embedding configuration must work together on the exact BERT-Tiny implementation. Some Transformer operations may remain floating point, and converting weights does not guarantee faster CPU inference. PyTorch's documented dynamic quantization support is module-specific, so compatibility must be established rather than assumed.

Required before submission: a successful small pilot, or a clearly documented technical uncertainty and fallback. Full training is not required.

Major 2The novelty claim is not sufficiently established.

Five references meet the minimum numerically, but the claim that prior studies do not examine this combination needs a broader literature check. The papers demonstrate related techniques; they do not prove the absence of earlier combined studies.

Major 3The evaluation scope is narrower than a general CV–job fit claim.

Your reported CV overlap is acknowledged, which is good. However, the benchmark cannot establish generalization to unseen candidates. The final report must consistently describe the intended unseen-job evaluation and avoid stronger claims.

Major 4The compression target needs an empirical justification.

The 0.40 size ratio is a legitimate, falsifiable target, but the 90\% embedding parameter share does not directly imply a 60\% reduction in total serialized size. Parameter count, storage dtype, quantization metadata, and unquantized components all matter.

The target is not inherently unreasonable. Present it as a hypothesis motivated by the architecture, not as a result implied by the embedding share.

\paragraph*{Minor issues}
\begin{itemize}
\item Teacher selection: Specify the validation metric used to select the teacher checkpoint and the exact stopping or checkpoint-selection rule.
\item Distillation reproducibility: State the temperature and loss-weight search procedure completely, including how the best configuration is chosen and whether the teacher's logits are computed in evaluation mode.
\item Three-seed uncertainty: Three seeds are acceptable for a small course project, but the mean and standard deviation are descriptive, not strong evidence of statistical significance. Your proposal already acknowledges this.
\item Prediction-agreement pilot criterion: At least 95\% agreement with M1 is a useful diagnostic, but it is not evidence that M2 preserves task performance. Keep the separate macro-F1 hypothesis.
\item Latency and memory: Specify the CPU model, operating system, software versions, and memory-measurement method. Distinguish model weight storage from runtime peak memory.
\item Fallback interpretation: A Linear-only quantization variant can be useful exploratory evidence, but it is a different M2. It must not be presented as a successful test of the original embedding-plus-Linear quantization hypothesis.
\item Label audit: Independently labeling 90 training pairs is useful. However, Cohen's kappa measures annotator agreement, not whether the dataset's labels reflect actual hiring outcomes.
\item Dataset provenance and license: The linked dataset is listed as containing 8,000 rows, and its repository exposes the pinned revision. The dataset page does not provide a dataset card, so provenance and labeling details remain unresolved.
\end{itemize}
Verify the license and permitted use of the underlying CV text before publication or redistribution.

\begin{itemize}
\item AI disclosure: This is appropriately transparent. Confirm that it meets Effat University's applicable disclosure requirements; I cannot determine course-specific compliance from the proposal alone.
\end{itemize}
\subsubsection*{5. References you must verify yourselves}
I checked the discoverable bibliographic records for the cited works. This is a preliminary metadata check, not a substitute for reading the original papers and verifying the reported numerical claims.

\begin{enumerate}
\item Devlin et al. (2019) — BERT
\end{enumerate}
Status: Bibliographic record identifiable.

Verify that your description of BERT supports the specific architecture and modeling choices you use, rather than treating the original BERT paper as evidence for every property of BERT-Tiny.

ACL Anthology: original paper

\begin{enumerate}
\item Turc et al. (2019 preprint; ICLR 2020)
\end{enumerate}
Status: Paper verified; numerical claim requires checking.

The work studies pretrained compact models and distillation. Verify the exact GLUE development score comparison (82.8 versus 84.4), including which student and evaluation setup the numbers describe. The work is also listed as an ICLR 2020 paper, so check the preferred publication citation.

arXiv: original paper

\begin{enumerate}
\item Hinton, Vinyals, and Dean (2015) — Distilling the Knowledge in a Neural Network
\end{enumerate}
Status: Original preprint identifiable; result needs verification.

Check the cited MNIST error counts (74 versus 146) against the original experiment and ensure the comparison is described with its proper model and training context.

arXiv: original paper

\begin{enumerate}
\item Zafrir et al. (2019) — Q8BERT
\end{enumerate}
Status: Bibliographic record identifiable; comparison requires checking.

Verify the MRPC F1 values (90.0, 89.56, and 87.88), their experimental settings, and whether the quoted results refer to the quantization procedures you describe. Q8BERT's quantization-aware training differs from your post-training approach.

arXiv: original paper

\begin{enumerate}
\item Gan, Zhang, and Mori (2024) — Application of LLM Agents in Recruitment
\end{enumerate}
Status: Journal record verified.

The authors, journal, volume, pages, year, and DOI are consistent with the publisher's record. Read the paper to verify that it supports the specific motivation attributed to it; its LLM-agent recruitment framework is not equivalent to your compact three-class classifier.

Publisher's article page

\paragraph*{Additional source verification}
Beyond the five cited papers, verify the following primary sources and technical claims:

\begin{itemize}
\item BERT-Tiny checkpoint: Confirm the provenance, configuration, vocabulary, and parameter count of \texttt{google/bert\_uncased\_L-2\_H-128\_A-2}. Do not assume that the original BERT paper establishes the exact 4.4M parameter count for your checkpoint.
\item PyTorch quantization documentation: Check the installed version's API behavior, supported modules, embedding configuration, and serialization procedure. The current documentation is a useful starting point, not proof of compatibility with your environment. PyTorch quantization recipe.
\item Dataset revision and split: Download the exact revision \texttt{08978e2}, independently run the audit, and verify every count after removing conflicting and duplicate pairs. Pinned dataset repository.
\item Prior combined methods: Search the references and related work of the five papers for studies combining distillation and post-training quantization. Until that search is complete, the gap is plausible but not conclusively established.
\end{itemize}
I have not independently reproduced the five papers' experimental results, your dataset audit, or the proposed quantization behavior. None of the proposed project's outcomes should be described as established findings.

\subsubsection*{6. Three questions your instructor may ask}
Question 1: What is genuinely your deep learning contribution if knowledge distillation already exists?

Be prepared to explain that the contribution is the controlled investigation of whether teacher-guided training improves a pretrained compact student on this task, using M0 as the matched hard-label baseline. Do not claim that the distillation algorithm itself is new.

Question 2: Why should quantizing the embeddings and Linear layers produce a 60\% storage reduction, and how will you prove it?

Explain that 0.40 is a target to test, not a guaranteed consequence of the embedding parameter share. Describe the pilot's conversion checks, saved-state-dict measurement, and planned comparison of M2 with M1 and M0.

Question 3: How can you trust your test results if the same CVs appear in training and test, and the dataset's labeling method is undocumented?

Explain the narrower unseen-job evaluation, the deduplication and overlap audit, and the independent 90-pair label audit. Acknowledge that these measures reduce some risks but do not establish generalization to new candidates or validate the labels as real hiring outcomes.

\subsubsection*{Final recommendation}
Readiness: READY TO SUBMIT by the stated score threshold, with important caveats.

Your strongest features are the explicit M0/M1/M2 definitions, the connected H1/H2 hypotheses, the frozen numerical targets, and the separation of validation from test evaluation. Those are substantial strengths for a formative proposal.

Before submitting, prioritize these three actions:

\begin{enumerate}
\item Validate the complete quantization and reload path on the actual checkpoint and intended PyTorch environment.
\item Substantiate and carefully scope the literature gap, including verification of all reported paper-specific numbers.
\item Make the data limitations and evaluation scope unmistakable, particularly shared CVs, undocumented labels, and the distinction between unseen jobs and unseen candidates.
\end{enumerate}
You do not need to train the full model to address these proposal-stage concerns. A small, documented pilot is sufficient to establish technical feasibility. If the pilot cannot verify the planned M2, or the literature review cannot support the claimed gap, resolve those issues before relying on the current readiness assessment.

}{}
\end{document}
